\subsection{部分等距线性映射}

定义 2. --- 设 E 与 F 为两个希尔伯特空间，\(u \in \mathcal{L}(E; F)\)。u 的核在 E 中的正交补称为 u 的初始子空间，u 的像在 F 中的闭包称为 u 的终子空间。从 E（相应地 F）到 u 的初始（相应地终）子空间上的正交投影算子称为 u 的初始（相应地终）正交投影算子。

设 \(\mathbf{P}\) 为 \(u\) 的初始子空间。由于 \(\mathbf{E}\) 是 \(\mathbf{P}\) 与 \(u\) 的核的直和，我们有 \(u(\mathbf{P}) = u(\mathbf{E})\)。

命题 4. --- (i) \(u^*\) 的初始（相应地终）子空间等于 \(u\) 的终（相应地初始）子空间。

\begin{enumerate}
\def\labelenumi{(\roman{enumi})}
\setcounter{enumi}{1}
\tightlist
\item
  假定 \(\mathbf{E} = \mathbf{F}\)。设 \(\mathbf{M}\) 为 \(\mathbf{E}\) 的闭向量子空间，\(\mathbf{M}^{\circ}\) 为其正交补。关系 \(u(\mathbf{M}) \subset \mathbf{M}\) 与 \(u^{*}(\mathbf{M}^{\circ}) \subset \mathbf{M}^{\circ}\) 等价。
\end{enumerate}

设 \(\mathrm{Q} = \overline{u(\mathrm{E})}\) 为 \(u\) 的终子空间。\(\mathrm{Q}^{\circ}\)（\(\mathrm{Q}\) 在 \(\mathrm{F}\) 中的正交补）由所有这样的向量 \(y\) 组成：\(\langle u(x)|y\rangle = 0\) 对一切 \(x\in \mathrm{E}\) 成立；这等价于：\(\langle x|u^{*}(y)\rangle = 0\) 对一切 \(x\in \mathrm{E}\) 成立，亦即 \(u^{*}(y) = 0\)。因此 \(\mathrm{Q}^{\circ} = \operatorname{Ker}u^{*}\)，而 \(\mathrm{Q}\) 是 \(u^{*}\) 的初始子空间。由于 \(u\) 是 \(u^{*}\) 的伴随，\(u^{*}\) 的终子空间也是 \(u\) 的初始子空间。这就证明了 (i)。

关系 \(u(\mathbf{M}) \subset \mathbf{M}\) 蕴含 \(u(\mathbf{M})\) 与 \(\mathbf{M}^{\circ}\) 正交，而关系 \(u^{*}(\mathbf{M}^{\circ}) \subset \mathbf{M}^{\circ}\) 蕴含 \(u^{*}(\mathbf{M}^{\circ})\) 与 \(\mathbf{M}\) 正交。但我们有 \(\langle u(x)|y\rangle = \overline{\langle u^{*}(y)|x\rangle}\)（对所有 \(x \in \mathbf{M}\) 与 \(y \in \mathbf{M}^{\circ}\)）；由此即得 (ii)。

我们注意到，利用 V, p.~41 的注 3，命题 4 可以由转置的一般性质（II, p.~51, 推论 2）推出。

定义 3. --- 设 E 与 F 为两个希尔伯特空间。称映射 \(u \in \mathcal{L}(\mathrm{E};\mathrm{F})\) 是部分等距的，如果对所有属于 u 的初始子空间的 x 有 \(\|u(x)\| = .\|x\|\)。

设 \(u \in \mathcal{L}(\mathrm{E};\mathrm{F})\)，\(\mathbf{N}\) 为其核，\(\mathbf{I}\) 为其像。说 \(u\) 是部分等距的，等于说线性映射 \(\tilde{u}: \mathrm{E/N} \to \mathrm{I}\)（由 \(u\) 导出的）是等距的（V, p.~13）。这时子空间 \(\mathbf{I}\)（\(\mathbf{F}\) 中的）是完备的，因而是闭的，并且是 \(u\) 的终子空间。于是，\(u\) 诱导出从 \(u\) 的初始子空间到其终子空间上的一个希尔伯特空间同构。

命题 5. --- 设 \(u \in \mathcal{L}(\mathrm{E};\mathrm{F})\)，\(\mathbf{P}\) 为其初始子空间，\(\mathbf{Q}\) 为其终子空间。以 \(p\)（相应地 \(q\)）表示 \(u\) 的初始（相应地终）正交投影算子。假定 \(u\) 是部分等距的。

\begin{enumerate}
\def\labelenumi{(\roman{enumi})}
\item
  映射 \(u^{*} \in \mathcal{L}(\mathbf{F};\mathbf{E})\) 是部分等距的，其初始子空间为 \(\mathbf{Q}\)，终子空间为 \(\mathbf{P}\)。这时从 \(\mathbf{P}\) 到 \(\mathbf{Q}\) 上的同构（由 \(u\) 诱导的）是从 \(\mathbf{Q}\) 到 \(\mathbf{P}\) 上的同构（由 \(u^{*}\) 诱导的）的逆。
\item
  我们有 \(u^{*}u = p\) 且 \(uu^{*} = q\)。
\end{enumerate}

根据命题 4 (i)，断言 (i) 是 (ii) 的推论。

我们现在证明 (ii)。由于 \(\mathbf{P}\) 包含 \(u^{*}\) 的像，映射 \(u^{*}u\) 把 \(\mathbf{E}\) 映到 \(\mathbf{P}\) 中。设 \(x \in \mathbf{E}\) 且 \(y \in \mathbf{P}\)，则

\[
\langle u ^ {*} u (x) | y \rangle = \langle u (x) | u (y) \rangle .
\]

若 x 属于 P，则由部分等距映射的定义有 \(\langle u(x)|u(y)\rangle = \langle x|y\rangle\)；若 x 属于 u 的核 N，则 \(u(x) = 0\)，从而 \(\langle u(x)|u(y)\rangle = 0\)，且由于 N 与 P 正交而有 \(\langle x|y\rangle = 0\)。因为 \(E = P \oplus N\)，所以在所有情形都有 \(\langle u^{*}u(x) - x|y\rangle = 0\)，于是 \(u^{*}u\) 是从 E 到 P 上的正交投影算子 p。\(uu^{*} = q\) 则由在上面的论证中把 u 与 \(u^{*}\) 互换得到。

命题 6. --- 对每个 \(u \in \mathcal{L}(\mathrm{E};\mathrm{F})\)，下列条件等价：

\begin{enumerate}
\def\labelenumi{(\roman{enumi})}
\item
  \(u\) 是部分等距的；
\item
  \(u^{*}\) 是部分等距的；
\item
  \(u^{*}u\) 是正交投影算子；
\item
  \(uu^{*}\) 是正交投影算子；
\item
  \(uu^{*}u = u;\)
\item
  \(u^{*}uu^{*} = u^{*}\)。
\end{enumerate}

由命题 5，(i) 与 (ii) 等价。

\begin{enumerate}
\def\labelenumi{(\roman{enumi})}
\item
  \(\Rightarrow\) (v)：假定 \(u\) 是部分等距的。于是由命题 5，\(u^{*}u\) 是 \(u\) 的初始正交投影算子。因此对每个 \(x \in \mathrm{E}\)，\(u^{*}u(x) - x\) 属于 \(u\) 的核，即 \(uu^{*}u(x) = u(x)\)。
\item
  \(\Rightarrow\) (iii)：假定 \(uu^{*}u = u\)，令 \(p = u^{*}u\)；则 \(p = p^{*}\) 且 \(p^2 = p\)。设 M（相应地 N）为 \(p\) 的像（相应地核）。对 \(x \in \mathbf{M}\) 与 \(y \in \mathbf{N}\)，我们有 \(\langle x|y\rangle = \langle p(x)|y\rangle = \langle x|p^{*}(y)\rangle = \langle x|p(y)\rangle = 0\)。由于 M 与 N 正交，\(p\) 是从 E 到 M 上的正交投影算子。
\item
  \(\Rightarrow\) (i)：假定 \(p = u^{*}u\) 是以 M 为像、以 N 为核的正交投影算子。
\end{enumerate}

对所有 \(x\in \mathbf{E}\)，我们有

\[
\left\| u (x) \right\| ^ {2} = \langle u ^ {*} u (x) | x \rangle = \langle p (x) | x \rangle .
\]

因此 \(u(x) = 0\)（对 \(x \in \mathbf{N}\)），且 \(\| u(x) \| = \| x \|\)（对 \(x \in \mathbf{M}\)），于是 \(u\) 是部分等距的，其核为 \(\mathbf{N}\)，初始子空间为 \(\mathbf{M}\)。

我们证明了 (i)、(iii) 与 (v) 的等价性。把 \(u\) 换成 \(u^*\)，即可推出 (ii)、(iv) 与 (vi) 的等价性。这就证明了命题 6。
% semantic-fragments: legacy-input-seam
\relax{}
